@@M0@@

\textbf{Remark 1.14.} Using the notion of a ``relative scheme'' (cf [3]), we can still give a meaning to 1.1 by assuming that @@M1@@ is a locally ringed topos. If we assume that the final object of the topos is quasi-compact, the proof given in the case of schemes carries over without difficulty, and the theorem therefore remains valid.

\sgaStableTarget{sga6.e06.sec.2}\section*{2. Computation of the @@M2@@ of a projective bundle: the case of perfect complexes}

The present section will not be used in the rest of the seminar.

In this section, the notation @@M3@@ will denote the Grothendieck group of perfect complexes on @@M4@@ (IV 2.2). We propose to prove the following theorem, analogous to 1.1:

\sgaStableTarget{sga6.e06.thm.2-1}\textbf{Theorem 2.1.} \uline{Let} @@M5@@ \uline{be a quasi-compact scheme}, @@M6@@ \uline{a locally free} @@M7@@\uline{-Module of rank} @@M8@@, @@M9@@ \uline{the associated projective bundle}; \uline{let} @@M10@@ \uline{be the class} \uline{in} @@M11@@ \uline{of the complex equal to} @@M12@@ \uline{in degree zero}, \uline{and zero in the other degrees}. \uline{Then} @@M13@@, \uline{endowed with its natural} @@M14@@-\uline{algebra structure, is} \uline{a free} @@M15@@-\uline{module with basis} @@M16@@.

\sgaStableTarget{sga6.e06.sec.2.2}2.2. Let us fix the notation: @@M17@@ will denote the symmetric algebra of @@M18@@; we shall write @@M19@@ for the category of @@M20@@-Modules, @@M21@@ for the category of graded @@M22@@-Modules with graded morphisms as morphisms, @@M23@@ for the full subcategory of @@M24@@ whose objects are the @@M25@@-Modules locally free of finite type as graded @@M26@@-Modules, and @@M27@@ for the category of @@M28@@-Modules. The categories @@M29@@, @@M30@@, and @@M31@@ are abelian.

For every @@M32@@, and every complex @@M33@@ in @@M34@@, we shall write @@M35@@ for the complex @@M36@@.

